%% toc: P2 9 9.1,9.2,9.3,9.4
\section{Proofs for the comparison with historical reconstructions}
\label{sec:annex-p02-w02-p02-001-historical-contrast}

Throughout, \(\mathfrak A\) is a structure with nonempty universe \(A\), as
in Definition~\ref{def:p02-syllogistic-semantics}; patterns, \(\models_{\mathrm B}\),
contradictories, \(\sigma\) and \(\tau_s\) are as in
Definitions~\ref{def:p02-w02-p02-001-patterns}
and~\ref{def:p02-w02-p02-001-selection}. For a pattern with subject \(p\)
and predicate \(q\), write \(S=p^{\mathfrak A}\) and \(P=q^{\mathfrak A}\).

\begin{annexitem}{annex:p02-w02-p02-001-agreement}{thm:p02-w02-p02-001-agreement}{Where the two evaluations agree}
\begin{proof}
By Definition~\ref{def:p02-syllogistic-semantics},
\(\mathfrak A\models\forall(p,q)\) iff \(S\subseteq P\);
\(\mathfrak A\models\forall(p,\bar q)\) iff \(S\subseteq A\setminus P\), that
is, iff \(S\cap P=\varnothing\);
\(\mathfrak A\models\exists(p,q)\) iff \(S\cap P\ne\varnothing\); and
\(\mathfrak A\models\exists(p,\bar q)\) iff \(S\cap(A\setminus P)\ne\varnothing\),
that is, iff \(S\setminus P\ne\varnothing\).

\emph{Item 1.} Let \(S\ne\varnothing\). Then \(\alpha(S,P)\) is
\(S\subseteq P\), because its conjunct \(S\ne\varnothing\) holds.
\(\iota(S,P)\) is \(S\cap P\ne\varnothing\) by definition, and
\(\varepsilon(S,P)\), its negation, is \(S\cap P=\varnothing\).
\(\omega(S,P)\) is the negation of \(S\subseteq P\), which holds iff some
member of \(S\) lies outside \(P\), that is, iff \(S\setminus P\ne\varnothing\).
Comparing with the four conditions of the first paragraph gives
\(\mathfrak A\models_{\mathrm B}\varphi\) iff \(\mathfrak A\models\sigma(\varphi)\)
for each of the four patterns. This is also the content of the first
sentence of Theorem~\ref{thm:p13-boger-existence-bridge}.

\emph{Item 2.} Let \(S=\varnothing\). Then \(\alpha(S,P)\) is false because
its conjunct \(S\ne\varnothing\) fails, and \(\omega(S,P)=\neg\alpha(S,P)\)
is true. \(S\cap P=\varnothing\), so \(\iota(S,P)\) is false and
\(\varepsilon(S,P)\) is true. Thus exactly the \(\mathsf e\)- and
\(\mathsf o\)-patterns hold under \(\models_{\mathrm B}\). On the other side,
\(\varnothing\subseteq P\) and \(\varnothing\cap P=\varnothing\) make
\(\forall(p,q)\) and \(\forall(p,\bar q)\) true, while
\(\varnothing\cap P\) and \(\varnothing\setminus P\) are empty, so
\(\exists(p,q)\) and \(\exists(p,\bar q)\) are false. Thus exactly the
\(\mathcal S\)-readings of the \(\mathsf a\)- and \(\mathsf e\)-patterns
hold.

\emph{Final sentence.} If \(S\ne\varnothing\), all four patterns agree with
their readings by item~1. If \(S=\varnothing\), item~2 shows that the
\(\mathsf e\)-pattern and its reading are both true, the
\(\mathsf i\)-pattern and its reading are both false, the
\(\mathsf a\)-pattern is false while its reading is true, and the
\(\mathsf o\)-pattern is true while its reading is false.
\end{proof}
\end{annexitem}

\begin{annexitem}{annex:p02-w02-p02-001-opposition}{thm:p02-w02-p02-001-opposition}{Opposition under both policies}
\begin{proof}
\emph{Item 1.} For \(\models_{\mathrm B}\), the contradictory pairs are
\(q\mathbin{\mathsf a}p,q\mathbin{\mathsf o}p\) and
\(q\mathbin{\mathsf e}p,q\mathbin{\mathsf i}p\), evaluated by \(\alpha,\omega\)
and \(\varepsilon,\iota\). By Definition~\ref{def:p13-boger-guarded-bridge},
\(\omega=\neg\alpha\) and \(\varepsilon=\neg\iota\), so exactly one member of
each pair holds, for empty and nonempty \(S\) alike. For the
\(\mathcal S\)-readings, \(\sigma(\varphi^{\mathrm c})=\overline{\sigma(\varphi)}\):
indeed \(\overline{\forall(p,q)}=\exists(p,\bar q)\) and
\(\overline{\forall(p,\bar q)}=\exists(p,q)\) by
Definition~\ref{def:p02-syllogistic-identifications}, and the formula bar is
an involution. Lemma~\ref{thm:p02-syllogistic-semantic-identifications}
states that exactly one of a formula and its complement is true.

\emph{Item 2.} Suppose \(\alpha(S,P)\) and \(\varepsilon(S,P)\) both hold.
The first gives \(S\ne\varnothing\) and \(S\subseteq P\); choose \(u\in S\).
Then \(u\in S\cap P\), so \(\iota(S,P)\) holds, contradicting
\(\varepsilon(S,P)=\neg\iota(S,P)\). For the \(\mathcal S\)-readings,
\(S\subseteq P\) and \(S\cap P=\varnothing\) together give
\(S=S\cap P=\varnothing\); conversely, \(S=\varnothing\) makes both
inclusions true (proof of Theorem~\ref{thm:p02-w02-p02-001-agreement}).

\emph{Item 3.} If \(\iota(S,P)\) fails, then \(S\cap P=\varnothing\). Either
\(S=\varnothing\), so \(\alpha(S,P)\) fails, or some \(u\in S\) exists, and
\(u\notin P\) shows \(S\not\subseteq P\), so again \(\alpha(S,P)\) fails.
In both cases \(\omega(S,P)=\neg\alpha(S,P)\) holds. For the
\(\mathcal S\)-readings, both existentials fail iff \(S\cap P=\varnothing\)
and \(S\setminus P=\varnothing\), iff \(S=(S\cap P)\cup(S\setminus P)=\varnothing\).

\emph{Item 4.} If \(\alpha(S,P)\), then \(S\ne\varnothing\) and
\(S\cap P=S\ne\varnothing\), so \(\iota(S,P)\). If \(\varepsilon(S,P)\)
holds and \(S=\varnothing\), then \(\alpha(S,P)\) fails and \(\omega(S,P)\)
holds. If \(\varepsilon(S,P)\) holds and \(u\in S\), then \(u\notin P\), so
\(S\not\subseteq P\), \(\alpha(S,P)\) fails and \(\omega(S,P)\) holds.
For the failures, let \(\mathfrak A\) have universe \(\{c\}\) and
\(p^{\mathfrak A}=q^{\mathfrak A}=\varnothing\). Then \(\forall(p,q)\) and
\(\forall(p,\bar q)\) are true, while \(\exists(p,q)\) and
\(\exists(p,\bar q)\) are false, by item~2 of
Theorem~\ref{thm:p02-w02-p02-001-agreement}.
\end{proof}
\end{annexitem}

\begin{annexitem}{annex:p02-w02-p02-001-o-inexpressible}{thm:p02-w02-p02-001-o-inexpressible}{The partial privative is inexpressible}
\begin{proof}
\emph{The four equivalences.} \(\alpha(S,P)\) is the conjunction of
\(S\ne\varnothing\), which is \(\mathfrak A\models\exists(p,p)\), and
\(S\subseteq P\), which is \(\mathfrak A\models\forall(p,q)\).
\(\varepsilon(S,P)\) is \(S\cap P=\varnothing\), which is
\(\mathfrak A\models\forall(p,\bar q)\), and \(\iota(S,P)\) is
\(\mathfrak A\models\exists(p,q)\). Finally \(\omega(S,P)=\neg\alpha(S,P)\)
holds iff \(S=\varnothing\) or \(S\not\subseteq P\). The first disjunct is
\(S\subseteq A\setminus S\), which is \(\mathfrak A\models\forall(p,\bar p)\);
the second is \(S\setminus P\ne\varnothing\), which is
\(\mathfrak A\models\exists(p,\bar q)\).

\emph{A transfer fact.} Call two elements \(e\) of \(\mathfrak A\) and
\(e'\) of \(\mathfrak B\) \emph{alike} when, for every atom \(o\),
\(e\in o^{\mathfrak A}\) iff \(e'\in o^{\mathfrak B}\). Then for every literal
\(l\), \(e\in l^{\mathfrak A}\) iff \(e'\in l^{\mathfrak B}\): for a positive
literal this is the definition, and for \(l=\bar o\) both sides say that
the element lies outside the extension of \(o\).

\emph{Four structures.} Let \(x,a,b\) be distinct objects. Every atom other
than \(p\) and \(q\) has empty extension in each of the following.
\(\mathfrak M_1\) has universe \(\{x\}\) and
\(p^{\mathfrak M_1}=q^{\mathfrak M_1}=\varnothing\).
\(\mathfrak U\) has universe \(\{x,a\}\) and
\(p^{\mathfrak U}=q^{\mathfrak U}=\{a\}\).
\(\mathfrak M_5\) has universe \(\{a,b\}\), \(p^{\mathfrak M_5}=\{a,b\}\)
and \(q^{\mathfrak M_5}=\{a\}\).
Then \(x\) in \(\mathfrak M_1\) and \(x\) in \(\mathfrak U\) are alike, and
\(a\) in \(\mathfrak U\) and \(a\) in \(\mathfrak M_5\) are alike.
In \(\mathfrak M_1\), \(p\) is empty, so \(\omega\) holds and
\(\mathfrak M_1\models_{\mathrm B}q\mathbin{\mathsf o}p\).
In \(\mathfrak U\), \(p^{\mathfrak U}=\{a\}\) is nonempty and included in
\(q^{\mathfrak U}\), so \(\alpha\) holds and
\(\mathfrak U\not\models_{\mathrm B}q\mathbin{\mathsf o}p\).
In \(\mathfrak M_5\), \(b\in p^{\mathfrak M_5}\setminus q^{\mathfrak M_5}\),
so \(\alpha\) fails and
\(\mathfrak M_5\models_{\mathrm B}q\mathbin{\mathsf o}p\).

\emph{The argument.} Suppose \(\Theta\subseteq\mathcal S^\dagger\) satisfies
\(\mathfrak A\models\Theta\) iff \(\mathfrak A\models_{\mathrm B}q\mathbin{\mathsf o}p\)
for every \(\mathfrak A\). Then \(\mathfrak M_1\models\Theta\),
\(\mathfrak U\not\models\Theta\) and \(\mathfrak M_5\models\Theta\). Choose
\(\theta\in\Theta\) false in \(\mathfrak U\).

If \(\theta=\exists(l,m)\), it is true in \(\mathfrak M_1\), so
\(x\in l^{\mathfrak M_1}\cap m^{\mathfrak M_1}\), since \(x\) is the only
element. By the transfer fact, \(x\in l^{\mathfrak U}\cap m^{\mathfrak U}\),
so \(\theta\) is true in \(\mathfrak U\), a contradiction.

Hence \(\theta=\forall(l,m)\), and some \(u\in l^{\mathfrak U}\setminus m^{\mathfrak U}\)
exists. If \(u=x\), the transfer fact gives
\(x\in l^{\mathfrak M_1}\setminus m^{\mathfrak M_1}\), so \(\theta\) is false
in \(\mathfrak M_1\), contradicting \(\mathfrak M_1\models\Theta\). So \(u=a\).
By the transfer fact, \(a\in l^{\mathfrak M_5}\setminus m^{\mathfrak M_5}\),
so \(\theta\) is false in \(\mathfrak M_5\), contradicting
\(\mathfrak M_5\models\Theta\). No such \(\Theta\) exists.
\end{proof}
\end{annexitem}

\begin{annexitem}{annex:p02-w02-p02-001-reduction}{thm:p02-w02-p02-001-reduction}{A reduction to \(\mathsf S\)}
\begin{proof}
\emph{Step 1: refutation form.} By item~1 of
Theorem~\ref{thm:p02-w02-p02-001-opposition}, a structure satisfies
\(\varphi^{\mathrm c}\) under \(\models_{\mathrm B}\) exactly when it does not
satisfy \(\varphi\). Hence \(\Gamma\models_{\mathrm B}\varphi\) holds iff no
structure satisfies every member of \(\Delta=\Gamma\cup\{\varphi^{\mathrm c}\}\)
under \(\models_{\mathrm B}\).

\emph{Step 2: selections.} Let \(\mathfrak A\) satisfy every member of
\(\Delta\) under \(\models_{\mathrm B}\). By the four equivalences of
Theorem~\ref{thm:p02-w02-p02-001-o-inexpressible}, \(\mathfrak A\) satisfies
\(\tau_s(\psi)\) for every \(\mathsf a\)-, \(\mathsf e\)- and
\(\mathsf i\)-pattern \(\psi\in\Delta\) and every selection \(s\), and for
each \(\mathsf o\)-pattern \(q\mathbin{\mathsf o}p\in\Delta\) it satisfies
at least one of \(\forall(p,\bar p)\), \(\exists(p,\bar q)\). Let \(s\) pick,
for each \(\mathsf o\)-pattern, a formula of the pair that \(\mathfrak A\)
satisfies. Then \(\mathfrak A\models\tau_s(\Delta)\). Conversely, if
\(\mathfrak A\models\tau_s(\Delta)\) for some selection \(s\), the same
equivalences give \(\mathfrak A\models_{\mathrm B}\psi\) for every
\(\psi\in\Delta\). Therefore \(\Delta\) is satisfiable under
\(\models_{\mathrm B}\) iff \(\tau_s(\Delta)\) is satisfiable for some \(s\).

\emph{Step 3: unsatisfiability in \(\mathsf S\).} The formula
\(\exists(r,\bar r)\) has the form \(\exists(p,l)\) and lies in
\(\mathcal S\); it is an absurdity and is false in every structure by
Lemma~\ref{thm:p02-syllogistic-semantic-identifications}. So a set
\(\Theta\subseteq\mathcal S\) is unsatisfiable iff
\(\Theta\models\exists(r,\bar r)\): if \(\Theta\) is unsatisfiable the
consequence holds vacuously, and if it holds, a structure satisfying
\(\Theta\) would satisfy a false formula. By
Theorem~\ref{thm:p02-syllogistic-s-completeness},
\(\Theta\models\exists(r,\bar r)\) iff
\(\Theta\vdash_{\mathsf S}\exists(r,\bar r)\).

Combining the three steps, \(\Gamma\models_{\mathrm B}\varphi\) iff
\(\tau_s(\Delta)\) is unsatisfiable for every \(s\), iff
\(\tau_s(\Delta)\vdash_{\mathsf S}\exists(r,\bar r)\) for every \(s\).
\(\Delta\) is finite, so it has finitely many \(\mathsf o\)-patterns and
finitely many selections, and each \(\tau_s(\Delta)\) is finite.
\end{proof}
\end{annexitem}

\begin{annexitem}{annex:p02-w02-p02-001-structural}{thm:p02-w02-p02-001-structural}{Structural properties of the modern derivability relations}
\begin{proof}
Derivations are the finite trees of
Definition~\ref{def:p02-syllogistic-derivation}: each leaf is a premise or
an instance of a rule without antecedents, and each other node follows by
a displayed rule from its immediate predecessors.

\emph{Item 1.} The one-node tree whose only node is \(\theta\) has a leaf
that is a premise in \(\Theta\); it derives \(\theta\).

\emph{Item 2.} A derivation from \(\Theta\) has every premise leaf in
\(\Theta\subseteq\Theta'\); the same tree is a derivation from \(\Theta'\).

\emph{Item 3.} Take a derivation of \(\theta\) from
\(\Theta\cup\{\theta_1,\ldots,\theta_n\}\). Replace every leaf labelled by
some \(\theta_i\) that is not in \(\Theta\) with a derivation of
\(\theta_i\) from \(\Theta\). Each inner node keeps its rule instance and
its immediate predecessors' formulas, and every leaf of the result is a
member of \(\Theta\) or a rule instance without antecedents. The result
is a derivation of \(\theta\) from \(\Theta\).

\emph{Item 4.} Rule X of either sheet has antecedents \(\psi,\bar\psi\) and
an arbitrary conclusion \(\varphi\) of the fragment. The tree with leaves
\(\psi,\bar\psi\) and root \(\varphi\) is a derivation from
\(\{\psi,\bar\psi\}\).

\emph{Item 5.} Rule T of \(\mathsf S\) has no antecedents and conclusion
\(\forall(p,p)\); rule T of \(\mathsf S^\dagger\), with \(l=p\), gives the
same formula. The one-node tree is a derivation from \(\varnothing\).

\emph{Item 6.} Both calculi are sound
(Theorems~\ref{thm:p02-syllogistic-sdagger-completeness}
and~\ref{thm:p02-syllogistic-s-completeness}). The structure with universe
\(\{c\}\) and \(p^{\mathfrak A}=q^{\mathfrak A}=\varnothing\) satisfies
\(\forall(p,q)\) and falsifies \(\exists(p,q)\)
(proof of Theorem~\ref{thm:p02-w02-p02-001-opposition}, item~4). A
derivation of \(\exists(p,q)\) from \(\{\forall(p,q)\}\) would make the
conclusion true in that structure.
\end{proof}
\end{annexitem}
